\section{Introduction}
\label{sec:intro}

Safety evaluations are the primary gatekeeper for deploying large language models (LLMs) in production.
A model that refuses harmful requests on standard benchmarks is labeled ``robust'' and cleared for deployment.
But what if this robustness is an artifact of evaluation scope?
Recent work on sleeper agents~\cite{hubinger2024sleeper}, emergent misalignment~\cite{betley2025thought}, and chain-of-thought (CoT) hijacking~\cite{chang2026unreal} reveals that models can maintain deceptive alignment --- behaving safely under scrutiny while retaining the capacity for harmful behavior in alternative contexts.
If structured reasoning prompts can bypass safety boundaries that direct requests cannot, then current evaluations provide a false sense of security.

{\bf Why CoT creates a new attack surface.}
Chain-of-thought prompting~\cite{wei2022chain, kojima2022large} improves model reasoning by eliciting step-by-step explanations.
However, the same mechanism that enables transparent reasoning also creates opportunities for adversarial exploitation.
The AdvChain framework~\cite{zhu2025advchain} identifies a ``Snowball Effect'' in which small reasoning errors amplify through CoT chains, and recent work shows that CoT reasoning can be systematically decoupled from model intent via backdoor injection~\cite{chang2026unreal}.
\citet{su2024enhancing} demonstrate that replacing standard affirmative targets with CoT triggers improves adversarial attack success.
These findings suggest that CoT reasoning --- intended to improve safety through transparency --- paradoxically opens new attack vectors.

{\bf The gap we address.}
Most prior work on hidden misalignment requires either training-time model access (\eg backdoor injection~\cite{hubinger2024sleeper, chang2026unreal}) or iterative optimization (\eg GCG~\cite{zou2023universal}, PAIR~\cite{chao2023pair}).
No systematic study exists of \emph{single-shot, black-box} adversarial CoT probing --- crafting inference-time prompts that exploit the reasoning process itself to surface misalignment in models that pass standard safety benchmarks.
The alignment faking literature~\cite{alignmentfaking2025} shows that models exhibit context-dependent compliance gaps, but this has not been systematically connected to adversarial CoT as a probing methodology.

{\bf Our approach.}
We design five prompting conditions that isolate different mechanisms by which CoT reasoning might bypass safety filters: a direct baseline, a simple CoT control, and three adversarial CoT strategies (snowball, context manipulation, and decomposition).
We evaluate each condition on 80 harmful behaviors from \advbench~\cite{zou2023universal} across \gptmini and \gptfour, using an LLM judge with a structured rubric and reporting attack success rates with Clopper-Pearson confidence intervals and McNemar's paired tests with Bonferroni correction.
We complement this with over-refusal measurement on \xstest~\cite{rottger2023xstest} to verify that our findings reflect genuine vulnerability rather than general behavioral degradation.

{\bf What we find.}
Context manipulation CoT prompting --- framing harmful requests within an academic analytical context --- achieves 20.0\% \asr on \gptmini compared to a 2.5\% direct-prompt baseline ($p=0.0001$, \cohenh$=0.61$), an 8$\times$ increase.
\gptfour shows the same directional pattern at lower magnitude (7.5\% vs.\ 0.0\%).
Other CoT strategies do not significantly differ from direct prompting, establishing that the adversarial \emph{mechanism} matters more than the CoT \emph{format}.

Our contributions are as follows:
\begin{itemize}[leftmargin=*,itemsep=0pt,topsep=0pt]
    \item We conduct the first systematic study of single-shot, black-box adversarial CoT probing as a methodology for uncovering hidden misalignment in deployed LLMs.
    \item We demonstrate that context manipulation CoT prompting reveals statistically significant vulnerabilities that standard direct-prompt evaluations miss, with an 8$\times$ increase in attack success rate on \gptmini ($p=0.0001$).
    \item We show that not all CoT strategies are equally effective --- only context manipulation produces significant gains --- establishing that context shifting, not reasoning format, drives the vulnerability.
    \item We provide evidence that standard safety benchmarks systematically underestimate model vulnerability, with implications for evaluation methodology and deployment decisions.
\end{itemize}
